% =====================================================================
%  Webseite Förderverein KF/JF Kreuzau -- Anleitung für den Vorstand
%  Vorlage: deb-support (debsupport.sty, siehe latex/README.md), Vereinskopf
%  Bilder: Bildschirmfotos einer Testinstallation mit Beispieldaten
%  Übersetzen mit:  pdflatex anleitung_vorstand_webseite.tex  (3x)
% =====================================================================
\documentclass[a4paper,fontsize=10pt]{scrartcl}
\usepackage[bericht]{debsupport}
\usepackage{float}
\setcounter{tocdepth}{1}

% Vereinskopf statt deb-support (Vorlage bleibt, nur die Absenderdaten wechseln)
\dssetup{firma=Förderverein KF/JF Kreuzau}
\makeatletter
\renewcommand*{\ds@kopf}{%
  \put(24.5,279.4){\makebox(0,0)[lb]{\dslogo[20]}}%
  \put(195.6,276.3){\makebox(0,0)[rb]{{\dsserif\fontsize{12}{14.45}\selectfont
     \begin{tabular}[b]{@{}r@{}}Kinder- und Jugendfeuerwehr\\Kreuzau\\
     foerderverein-kf-jf-kreuzau.de\end{tabular}}}}%
  \put(17.8,271.85){\rule{179.8mm}{0.4pt}}}
\renewcommand*{\dsfolgekopf}[1]{%
  \parbox[b]{\textwidth}{\dslogo[11]\hfill #1\\[-0.4ex]%
    \rule{\textwidth}{0.4pt}}}
\makeatother
\renewcommand*{\dsdokumentart}{Anleitung}
\renewcommand*{\dskurztitel}{Webseite -- Anleitung für den Vorstand}
\renewcommand*{\dsbeleg}{}
\setlength{\emergencystretch}{3em}
\graphicspath{{bilder/}}

% Bildschirmfoto mit dünnem Rahmen:  \foto[Breite]{Datei}{Bildunterschrift}
\newcommand{\foto}[3][0.72\textwidth]{%
  \begin{figure}[H]
    \centering
    \fcolorbox{dsgrau!60}{white}{\includegraphics[width=#1]{#2}}
    \caption{#3}
  \end{figure}}
\newcommand*{\knopf}[1]{\textbf{\guillemotright #1\guillemotleft}}
\newcommand*{\stufe}[1]{\fbox{\small #1}}
\setcapindent{0pt}\addtokomafont{caption}{\small}\addtokomafont{captionlabel}{\bfseries}

\begin{document}

\dstitel{Die Webseite des Fördervereins}
        {Anleitung für den Vorstand: Seiten und Anhänge pflegen, Versionsverlauf,
         Newsletter, Benutzer verwalten}
        {Webseite: & \url{https://foerderverein-kf-jf-kreuzau.de}\\
         Für:      & Vorstandsmitglieder (Stufe mit Bearbeitungsrecht)\\
         Stand:    & \today}

\begin{dshinweis}[title=Kurz und knapp]
\begin{enumerate}
  \item Anmelden mit Passwort \emph{und} sechsstelligem Code aus der
        Authenticator-App; die Zwei-Faktor-Anmeldung ist für den Vorstand Pflicht.
  \item \knopf{+ Neue Seite} auf der Startseite: Titel, Datum, \emph{Sichtbar für},
        Text, Dateien. \knopf{Speichern}.
  \item Unter jeder Seite: \emph{Bearbeiten · Versionsverlauf · Als Newsletter
        versenden · Löschen}.
  \item Personen anlegen und einladen: \emph{Verwaltung} $\rightarrow$
        \emph{Benutzer} $\rightarrow$ \knopf{Benutzer hinzufügen}.
\end{enumerate}
\end{dshinweis}

\tableofcontents

% ------------------------------------------------------------------
\section{Überblick}

Diese Anleitung ergänzt die \emph{Anleitung für Mitglieder}, die Anmelden,
Lesen, \emph{Mein Konto} und \emph{Passwort vergessen} beschreibt. Hier geht es
um das, was nur der Vorstand kann:

\begin{itemize}
  \item Seiten anlegen, ändern und löschen, Dateien und Bilder anhängen,
  \item frühere Fassungen einer Seite ansehen und wiederherstellen,
  \item Seiten als Newsletter per E-Mail verschicken,
  \item Personen anlegen, einladen, einer Stufe zuordnen und deaktivieren.
\end{itemize}

Nach der Anmeldung stehen oben rechts zusätzlich \emph{Newsletter} und
\emph{Verwaltung}; auf der Startseite gibt es den Knopf \knopf{+ Neue Seite}.
Der Vorstand sieht alle Seiten, auch die mit \stufe{Vorstand}.

\foto[0.6\textwidth]{v03b_kopfzeile_vorstand}{Kopfzeile für den Vorstand}

\subsection*{Stufen}

Jede Seite hat eine Stufe (\emph{Sichtbar für}). Sie ist sichtbar für diese Stufe
und alle höheren:

\begin{center}
\begin{tabularx}{0.9\textwidth}{@{}lX@{}}
\toprule
\textbf{Sichtbar für} & \textbf{wer sieht die Seite} \\
\midrule
Öffentlich (ohne Anmeldung) & alle, auch Suchmaschinen \\
Fördermitglieder            & Fördermitglieder, Mitglieder, Vorstand \\
Mitglieder                  & Mitglieder, Vorstand \\
Vorstand                    & nur der Vorstand \\
\bottomrule
\end{tabularx}
\end{center}

Anhänge und Bilder haben immer dieselbe Sichtbarkeit wie ihre Seite.

% ------------------------------------------------------------------
\section{Anmelden mit Zwei-Faktor}

Wer Seiten bearbeiten darf, muss sich zusätzlich mit einem Code aus einer
Authenticator-App auf dem Smartphone anmelden, etwa \emph{Aegis}, \emph{FreeOTP},
\emph{Google} oder \emph{Microsoft Authenticator}. So reicht ein ausgespähtes
Passwort allein nicht, um Inhalte zu ändern oder Mails an alle Mitglieder zu
verschicken.

\subsection*{Einrichten (einmalig)}

Nach der ersten Anmeldung führt jede Seite zur Einrichtung:

\begin{enumerate}
  \item In der App ein neues Konto hinzufügen und den QR-Code scannen.
  \item Den sechsstelligen Code aus der App bei \emph{Authentifizierungscode}
        eintragen und \knopf{Aktivieren} klicken.
\end{enumerate}

\foto[0.62\textwidth]{v01_zwei_faktor_pflicht}{Einrichtung der Zwei-Faktor-Anmeldung
  beim ersten Anmelden}

\subsection*{Anmelden}

Nach Benutzername und Passwort fragt die Webseite nach dem aktuellen Code aus
der App. Der Code wechselt alle 30~Sekunden.

\foto[0.62\textwidth]{v02_code_eingeben}{Code aus der App eingeben}

\begin{dshinweis}[title=Für den Notfall vorsorgen]
Unter \emph{Mein Konto} $\rightarrow$ \emph{Zwei-Faktor-Anmeldung}
$\rightarrow$ \knopf{verwalten} \textbf{Wiederherstellungscodes} erzeugen,
ausdrucken und sicher aufbewahren. Damit kommen Sie auch ohne Smartphone in
Ihr Konto. Ohne Codes kann nur die Betreuung der Webseite die
Zwei-Faktor-Anmeldung zurücksetzen.
\end{dshinweis}

% ------------------------------------------------------------------
\section{Eine neue Seite anlegen}

Auf der Startseite oben auf \knopf{+ Neue Seite} klicken.

\foto[0.55\textwidth]{v04_neue_seite_leer}{Das Formular für eine neue Seite}

\begin{description}[style=nextline,leftmargin=0pt,itemsep=0.4\baselineskip]
  \item[Titel] erscheint als Überschrift und in der Übersicht der Startseite.
  \item[Datum] bestimmt die Reihenfolge auf der Startseite (neueste oben).
    Vorgabe ist heute; für Termine besser das Datum der Veröffentlichung
    lassen und den Termin im Text nennen.
  \item[Sichtbar für] ist voreingestellt auf \stufe{Mitglieder}. Für alle sichtbar
    wird eine Seite nur, wenn Sie bewusst \emph{Öffentlich (ohne Anmeldung)}
    wählen.
  \item[Inhalt] schreiben Sie im Editor wie in einer Textverarbeitung.
  \item[Dateien und Bilder hinzufügen] öffnet die Dateiauswahl; mehrere
    Dateien auf einmal sind möglich (je höchstens 20~MB).
\end{description}

\subsection{Der Editor}

\foto[0.75\textwidth]{v05b_werkzeugleiste}{Werkzeugleiste des Editors}

Von links: \emph{Rückgängig} und \emph{Wiederholen}; Absatzformat
(\emph{Absatz}, \emph{Überschrift}, \emph{Unterüberschrift}, \emph{Zitat});
\textbf{fett}, \emph{kursiv}, \underline{unterstrichen}; Aufzählung und
nummerierte Liste; Link einfügen; Tabelle; Formatierung entfernen.

\begin{itemize}
  \item Text aus Word oder einer E-Mail kann eingefügt werden. Schriftarten,
        Farben und Größen werden dabei entfernt, damit alle Seiten einheitlich
        aussehen.
  \item Bilder kommen nicht in den Text, sondern als Datei in den Bereich
        \emph{Anhänge} (siehe unten).
  \item Für Links die Adresse markieren lassen oder ein Wort markieren und
        auf das Kettensymbol klicken.
\end{itemize}

\subsection{Dateien und Bilder}

\begin{itemize}
  \item \textbf{Dateien} (PDF, Word, Excel …) erscheinen unter der Seite als
        Liste \emph{Anhänge}.
  \item \textbf{Bilder} (JPG, PNG, WebP) erscheinen auf der Seite: das erste
        als großes \emph{Titelbild} über dem Text (und als Vorschaubild auf der
        Startseite), weitere als \emph{Galerie} darunter. Große Fotos werden
        automatisch verkleinert, Standortangaben (GPS) entfernt.
\end{itemize}

\foto[0.55\textwidth]{v05_neue_seite_ausgefuellt}{Ausgefülltes Formular:
  öffentliche Seite mit zwei Bildern und einem PDF}

Nach \knopf{Speichern} sehen Sie die fertige Seite. Darunter stehen die
Vorstandsfunktionen \emph{Bearbeiten · Versionsverlauf · Als Newsletter
versenden · Löschen}.

\foto[0.5\textwidth]{v06_seite_gespeichert}{Die gespeicherte Seite mit
  Titelbild, Galerie und Anhang}

\begin{dshinweis}[title=Fotos von Kindern und Jugendlichen]
Fotos, auf denen Personen erkennbar sind, nur mit deren Einverständnis
veröffentlichen, bei Kindern und Jugendlichen mit dem der Eltern. Im Zweifel
die Seite nicht öffentlich, sondern für \stufe{Mitglieder} anlegen oder ein
anderes Bild wählen.
\end{dshinweis}

% ------------------------------------------------------------------
\section{Eine Seite ändern}

Unter der Seite auf \emph{Bearbeiten} klicken. Es öffnet sich dasselbe Formular
wie beim Anlegen. Im Bereich \emph{Anhänge} stehen jetzt die vorhandenen
Dateien:

\foto[0.8\textwidth]{v07_anhaenge}{Anhänge einer Seite bearbeiten}

\begin{description}[style=nextline,leftmargin=0pt,itemsep=0.4\baselineskip]
  \item[Bezeichnung / Bildunterschrift] Bei Dateien der Text, der statt des
    Dateinamens in der Liste steht (z.\,B.\ \emph{Flyer zur Aktion (PDF)}); bei
    Bildern die Bildunterschrift.
  \item[Titelbild] Haken beim Bild, das oben groß stehen soll. Ohne Haken wird
    das erste Bild Titelbild. Nur ein Titelbild ist möglich.
  \item[Entfernen] Haken setzen, dann \knopf{Speichern}: die Datei wird gelöscht.
\end{description}

Neue Dateien fügen Sie wie beim Anlegen hinzu. Mit \emph{Abbrechen} verlassen Sie
das Formular ohne Änderungen.

% ------------------------------------------------------------------
\section{Versionsverlauf}

Jedes Speichern legt eine neue Fassung an. Unter der Seite öffnet
\emph{Versionsverlauf} die Liste aller Fassungen mit Zeitpunkt und Person.

\foto[0.62\textwidth]{v08_verlauf}{Versionsverlauf einer Seite}

\emph{Ansehen und vergleichen} zeigt eine frühere Fassung: oben die
Unterschiede zur aktuellen Fassung (grün: nur aktuell, rot: nur früher,
gelb: geändert), darunter die Seite, wie sie damals aussah.

\foto[0.55\textwidth]{v09_fassung}{Frühere Fassung: Uhrzeit von 14 auf 15 Uhr geändert}

\knopf{Diese Fassung wiederherstellen} übernimmt Titel, Datum, Sichtbarkeit und
Text der früheren Fassung. Die Anhänge bleiben, wie sie sind. Die bisherige
Fassung geht dabei nicht verloren, sie bleibt im Verlauf.

% ------------------------------------------------------------------
\section{Eine Seite löschen}

Unter der Seite auf \emph{Löschen} klicken und mit \knopf{Endgültig löschen}
bestätigen. Dabei werden auch alle Anhänge der Seite gelöscht.

\foto[0.62\textwidth]{v10_loeschen}{Rückfrage vor dem Löschen}

\begin{dshinweis}[title=Lieber verstecken als löschen]
Eine gelöschte Seite lässt sich auf der Webseite nicht zurückholen, nur die
Betreuung kann sie aus der Datensicherung wiederherstellen. Soll eine Seite nur
nicht mehr zu sehen sein, stellen Sie stattdessen \emph{Sichtbar für} auf
\stufe{Vorstand}.
\end{dshinweis}

% ------------------------------------------------------------------
\section{Newsletter}

Jede Seite kann per E-Mail an die Personen verschickt werden, die sie sehen
dürfen, beliebig oft. Typisch: die Einladung zur Mitgliederversammlung.

\subsection{Versand vorbereiten}

Unter der Seite auf \emph{Als Newsletter versenden} klicken.

\foto[0.58\textwidth]{v11_newsletter_formular}{Newsletter vorbereiten}

\begin{description}[style=nextline,leftmargin=0pt,itemsep=0.4\baselineskip]
  \item[An] Zur Wahl stehen die Stufe der Seite und alle höheren. Mit den Haken
    können Sie den Kreis einschränken, z.\,B.\ eine Information nur an den
    Vorstand. Bei öffentlichen Seiten gibt es zusätzlich \emph{Konten ohne
    Stufe}. Höhere Stufen erhalten die E-Mail zuerst.
  \item[Inhalt] \emph{Seite in der E-Mail}: der ganze Text steht in der Mail.
    \emph{Nur Vorschau und Link}: die ersten Zeilen und ein Knopf
    \knopf{Weiterlesen}. Bei Seiten für Mitglieder ist dann zum Lesen eine
    Anmeldung nötig; der Text verlässt den geschützten Bereich nicht.
  \item[Ankündigung] Ein kurzer, optionaler Text, der über der Seite steht.
  \item[Anhänge mitsenden] Dateien (keine Bilder) gehen als Anhang mit, solange
    sie zusammen höchstens 8~MB groß sind; sonst stehen Links in der Mail.
\end{description}

Unter dem Formular steht, wie viele Personen den Newsletter erhalten und wer
\textbf{keine E-Mail-Adresse} hat. Diese Personen müssen Sie auf anderem Weg
benachrichtigen, z.\,B.\ per Brief.

\subsection{Probe, Pause, Versand}

Nach \knopf{Versand starten} geht die E-Mail zuerst \textbf{nur an Sie selbst}.
Danach folgt eine Pause von 5~Minuten. Prüfen Sie in dieser Zeit die Mail in
Ihrem Postfach.

\foto[0.58\textwidth]{v12_newsletter_pause}{Die Probe ist verschickt, der Versand an alle wartet}

\begin{itemize}
  \item Alles in Ordnung: \knopf{Probe ist in Ordnung -- jetzt an alle senden}
        (oder die Pause abwarten).
  \item Fehler entdeckt: \knopf{Versand abbrechen}, die Seite korrigieren und
        den Versand neu starten.
\end{itemize}

Danach geht jede E-Mail einzeln hinaus. Die Seite aktualisiert sich von selbst
und zeigt für jede Person \emph{gesendet}, \emph{wartet}, \emph{Fehler} (mit
Meldung) oder \emph{keine E-Mail-Adresse}. Abbrechen ist bis zur letzten Mail
möglich; bereits verschickte Mails lassen sich nicht zurückholen.

\foto[0.58\textwidth]{v13_newsletter_fertig}{Abgeschlossener Versand mit Protokoll}

\subsection{Absender und Antworten}

Absender ist die Person, die den Newsletter verschickt: mit ihrer Vereinsadresse,
falls vorhanden, sonst \emph{Name (Förderverein KF/JF Kreuzau)} mit der Adresse
der Webseite. Antworten gehen an diese Person und an die Vorstandsadresse.

\subsection{Protokoll}

\emph{Newsletter} in der Kopfzeile listet alle Versände mit Zeitpunkt, Seite,
Absender und Status. Ein Klick auf den Zeitpunkt öffnet das Protokoll mit allen
Empfängern. Es dient z.\,B.\ als Nachweis, wann und an wen die Einladung zur
Mitgliederversammlung verschickt wurde (Satzung §\,7).

\foto[0.6\textwidth]{v14_newsletter_uebersicht}{Übersicht aller Versände}

% ------------------------------------------------------------------
\section{Benutzer verwalten}

Personen werden in der \emph{Verwaltung} gepflegt (Link oben rechts). Die
Verwaltung sieht etwas anders aus als die Webseite.

\foto[0.8\textwidth]{v15_verwaltung}{Startseite der Verwaltung}

\subsection{Benutzerliste}

\emph{Benutzer} zeigt alle Konten mit E-Mail-Adresse, Vereinsadresse, Stufe und
letzter Anmeldung. Rechts filtern Sie nach Stufe oder nach aktiven Konten, oben
suchen Sie nach Name oder Adresse.

\foto[0.95\textwidth]{v16_benutzerliste}{Benutzerliste}

\subsection{Person anlegen und einladen}

\knopf{Benutzer hinzufügen} (oben rechts in der Liste):

\foto[0.8\textwidth]{v17_benutzer_anlegen}{Neue Person anlegen}

\begin{description}[style=nextline,leftmargin=0pt,itemsep=0.4\baselineskip]
  \item[Benutzername, Vorname, Nachname] Benutzername klein und ohne Umlaute,
    z.\,B.\ \emph{vorname.nachname}. Anmelden kann man sich auch mit der
    E-Mail-Adresse.
  \item[E-Mail-Adresse] die private (externe) Adresse. Pflicht; an sie gehen
    Einladung und \emph{Passwort vergessen}. Jede Adresse kann nur zu einem
    Konto gehören.
  \item[Stufe] meist genau eine. \stufe{Vorstand} gibt Bearbeitungsrechte und
    macht die Zwei-Faktor-Anmeldung zur Pflicht.
  \item[Einladung sofort per E-Mail senden] ist angehakt: Nach dem Speichern
    erhält die Person eine Einladung mit einem Link, über den sie ihr Passwort
    selbst festlegt (7~Tage gültig).
  \item[Postfach und Vereinsmails] nur ausfüllen, wenn die Person eine
    Vereinsadresse (Postfach) hat; sonst leer lassen.
\end{description}

Mit \knopf{Sichern} wird das Konto angelegt und die Einladung verschickt.
Oben erscheint die Bestätigung \emph{Einladung an … gesendet}.

\subsection{Einladung erneut senden}

Ist der Link abgelaufen oder die Mail nicht angekommen: in der Benutzerliste
die Personen ankreuzen, bei \emph{Aktion} \emph{Einladung (erneut) senden}
wählen und \knopf{Ausführen} klicken.

\foto[0.5\textwidth]{v20_aktion_einladen}{Einladung an ausgewählte Personen senden}

\foto[0.6\textwidth]{v21_einladung_meldungen}{Rückmeldung: eine Einladung
  gesendet, eine Person ohne E-Mail-Adresse}

Der Vorstand kann keine Passwörter setzen oder ansehen. Hat jemand sein
Passwort vergessen, hilft \emph{Passwort vergessen} auf der Anmeldeseite oder
eine neue Einladung.

\subsection{Person ändern}

Ein Klick auf den Benutzernamen öffnet das Konto. Hier ändern Sie Name,
E-Mail-Adresse, Stufe und die Angaben zum Postfach.

\foto[0.75\textwidth]{v19_benutzer_bearbeiten}{Konto einer Person mit Vereinsadresse}

\begin{itemize}
  \item \textbf{Stufe wechseln}, etwa vom Fördermitglied zum Mitglied: Haken
        umsetzen und \knopf{Sichern}.
  \item \textbf{Ausgeschieden}: den Haken bei \emph{Aktiv} entfernen statt das
        Konto zu löschen. Die Person kann sich nicht mehr anmelden und erhält
        keine Newsletter; Verlauf und Protokolle bleiben vollständig.
  \item \textbf{Löschen} (rot, unten rechts) nur, wenn ein Konto versehentlich
        angelegt wurde oder die Person die Löschung ihrer Daten verlangt.
\end{itemize}

% ------------------------------------------------------------------
\section{Häufige Fragen}

\begin{description}[style=nextline,leftmargin=0pt,itemsep=0.6\baselineskip]
  \item[Ein Mitglied sieht eine Seite nicht.]
    \emph{Sichtbar für} der Seite und die Stufe der Person (Verwaltung)
    vergleichen. Ist die Person angemeldet?
  \item[Eine Seite war versehentlich öffentlich.]
    \emph{Sichtbar für} ändern und speichern. Ein bereits verschickter
    Newsletter lässt sich nicht zurückholen.
  \item[Ich habe etwas falsch geändert.]
    \emph{Versionsverlauf} öffnen und die vorige Fassung wiederherstellen.
  \item[Die Einladung kam nicht an.]
    E-Mail-Adresse in der Verwaltung prüfen, Person in den Spam-Ordner schauen
    lassen, Einladung erneut senden.
  \item[Im Newsletter-Protokoll steht „Fehler“.]
    Die Meldung daneben nennt den Grund, oft eine falsche oder volle Adresse.
    Adresse in der Verwaltung korrigieren und die Person auf anderem Weg
    informieren, z.\,B.\ mit einer kurzen E-Mail von Hand.
  \item[Die Statusseite meldet „Der Versanddienst meldet sich nicht“.]
    Dann verschickt der Server gerade keine Mails. Bitte die Betreuung
    informieren.
  \item[Smartphone mit der Authenticator-App verloren.]
    Mit einem Wiederherstellungscode anmelden und die App neu einrichten;
    ohne Codes die Betreuung bitten, die Zwei-Faktor-Anmeldung zurückzusetzen.
\end{description}

% ------------------------------------------------------------------
\section{Was die Betreuung der Webseite übernimmt}

Einige Einstellungen sind bewusst nur für Administratoren erreichbar:

\begin{itemize}
  \item \textbf{Stufen}: anlegen, umbenennen, festlegen, welche Stufe bearbeiten darf,
  \item \textbf{Einstellungen der Webseite}: Name, Slogan, Fußzeile,
        Antwortadresse für Newsletter,
  \item Server, Updates, Datensicherung, Wiederherstellen gelöschter Seiten,
  \item Zurücksetzen der Zwei-Faktor-Anmeldung.
\end{itemize}

{\small\color{dsgrau} Die Bildschirmfotos stammen aus einer Testinstallation
mit Beispieldaten; Personen, Seiten und Bilder sind erfunden. Die Beschriftung
der Dateiauswahl (\emph{Choose Files}) erscheint in Ihrem Browser auf Deutsch.}

\end{document}
